\documentclass[10pt]{article}

\usepackage[utf8]{inputenc}

\usepackage[T1]{fontenc}

\usepackage{amsmath}

\usepackage{amsfonts}

\usepackage{amssymb}

\usepackage[version=4]{mhchem}

\usepackage{stmaryrd}

\usepackage{graphicx}

\usepackage[export]{adjustbox}

\graphicspath{ {./images/} }



\begin{document}

$R I^{*}$-рруппа - обладает разрешимой нормальной П. с., вполне упорядоченной по возрастанию;



$\overline{R I}$-группа - всякую нормальную П. с. этой группы можно уплотнить до разрешимой нормальной;



$Z$-группа - обладает центральной П. с.;



$Z A$-группа - обладает центральной П. с., вполне упорядоченной по возрастанию;



$Z D$-группа - обладает центральной П. с., вполне упорядоченної по убыванию;



$\bar{Z}$-группа - всякую нормальную П. с. такой групшы можно уплотнить до центральной;



$\tilde{N}$-группа - через всякую подгруппу такой групшы проходит субнормальная П. с.;



$N$-группа - через всякую подгруппу такой группы проходит субнормальная П. с., вполне унорядоченная по возрастанию.



Частный случай П. с.- подгрупп рлды.



Лum.: [1] К у р о ш А. Г., Теория групп, 3 изд., М., 1967; [2] ч е р н и к о в С. Н., Группы с заданными свойствами системы подгрупп, М., 1980 ; [3] К а р г а п о ло в М. И., М е р 3л я к о в Ю. И, Оссновы теории групп, 3 изд., М., 1982 .



ПОДГРУППА - подмножество $H$ группы $G$, само являющееся груіпой относительно операции, определяющей $G$. Подмножество $H$ группы $G$ является ее подгруппой тогда и только тогда, когда: (1) $H$ содержит произведение любых двух элементов из $H$, (2) $H$ содержит вместе со всяким своим әлемевтом $h$ обратный к нему элемен'г $h^{-1}$. В случае конечных и, вообще, периодич. групп проверка условия (2) является излишней.



Подмножество группы $G$, состоящее из одного әлемента 1, будет, очевидно, подгруппой, и эта П. наз. е д и в и ч н ой П. группы $G$ и обозначается обычно символом $E$. Сама $G$ также является своей П. Всякая П., отличная от всей группы, наз. и с т и в в о й П. этой группы. Истинная П. нек-рой бесконечной группы может быть изоморфна самой группе. Сама группа $G$ и подгруппа $E$ наз. н е с о б с т в е н н ы м и П. группы $G$, все остальные - с о б с т в е н в ы м и.



Теоретико-множественное юересечение любых двух (и любого множества) П. группы $G$ является П. группы $G$. Пересечение всех П. группы $G$, содержащих все элементы нек-рого непустого множества $M$, наз. п о дг р уппой, порожденной множество м $M$, и обозначается символом $\{M\}$. Если $M$ состоит из одного элемента $a$, то $\{a\}$ наз. ц и к л и ч е с к о й П. әлемента $a$. Группа, совпадающая с одной из своих циклических П., наз. чиклической әруппой.



Теоретико-множественное объединение П., вообще говоря, не обязано являться П. О б ъ е ди н е ни е м п о д г р у п п $H_{i}, i \in I$, наз. П., порожденная объединением множеств $H_{i}$.



Произведение подмножеств $S_{1}$ и $S_{2}$ группы $G$ есть множество, состоящее из всевозможных (различных) произведевий $s_{1} s_{2}$, где $s_{1} \subset S_{2}, s_{2} \in S_{2}$. П роизведение подгрупп $\mathrm{H}_{1} \mathrm{H}_{2}$ есть П. тогда и только тогда, когда $\mathrm{H}_{1} \mathrm{H}_{2}=\mathrm{H}_{2} \mathrm{H}_{1}$, и в этом случае произведение $\mathrm{H}_{1} \mathrm{H}_{2}$ совпадает с объединением подгрупп $H_{1}$ и $\mathrm{H}_{2}$.



Гомоморфный образ П.- подгруппа. Если группа $G_{1}$ изоморфна век-рой подгруппе $H$ группы $G$, то говорят, что группа $G_{1}$ может быть вложена в группу $G$. Если даны две группы и каждая из вих изоморфна век-рой истинной П. другой, то отсюда еще не следует изоморфизм самих этих групі.



О. А. Кванова.



ПОДГРУППЫ ИНДЕКС смежных классов в каждом из разложений группы $G$ по этой подгруппе $H$ (в бесконечном случае - мощность множества этих классов). Если число смежных классов конечно, то $H$ наз. п о д г р у п п о й к он е ч в о г о и н д е к с а в $G$. Пересечение конечного числа подгрупп конечного индекса само имеет конеч- ный индекс (т е о р е м а П у а н к а p e). Индекс подгруппы $H$ в группе $G$ обычно обозначается $|G: H|$. Произведение порядка подгруппы $H$ на ее индекс $|G: H|$ равно порядку группы $G$ (т е о р е м а Л а гр а н ж а). Это соотношение имеет место как для конечной группы $G$, так и в случае бесконечной $G-$ для соответствующих мощностей.



\includegraphics[max width=\textwidth, center]{2023_07_08_d3ead3349a4727a6aeefg-1(1)}

Ю. И., Основы теории групI, 3 изд, M., 1982;

А. Г., Теория групі, 3 изд., М., 1967. иванова.



ПӧДЕРА, п о д э р а, кривой $l$ относительно точки $O-$ множество оснований перпендикуляров, опущенных из точки $O$ на касательные к кривой l. Напр., улитка Паскаля П. окружности относительно точки $O$ (см. рис.). П. плоской линии $x=x(t), \quad y=y(t)$ относительно начала координат:



\begin{center}

\includegraphics[max width=\textwidth]{2023_07_08_d3ead3349a4727a6aeefg-1}

\end{center}



\[

X=x-x^{\prime} \frac{x x^{\prime}+y y^{\prime}}{x^{\prime 2}+y^{\prime 2}}, Y=y-y^{\prime} \frac{x x^{\prime}+y y^{\prime}}{x^{\prime 2}+y^{\prime 2}} .

\]



Уравнение П. простравственной линии $x=x(t), y=$ $=y(t), z=z(t)$ относительно начала координат:



\[

\begin{gathered}

X=x-x^{\prime} \frac{x x^{\prime}+y y^{\prime}+z z^{\prime}}{x^{\prime 2}+y^{\prime 2}+z^{\prime 2}} ; Y=y-y^{\prime} \frac{x x^{\prime}+y y^{\prime}+z z^{\prime}}{x^{\prime 2}+y^{\prime 2}+z^{\prime 2}} ; \\

Z=z-z^{\prime} \frac{x x^{\prime}+y y^{\prime}+z z^{\prime}}{x^{\prime 2}+y^{\prime 2}+z^{\prime 2}} .

\end{gathered}

\]



А в т и п о д е р о й линии $l$ относительно точки $O$ наз. линия, П. к-рой относительно точки $O$ есть линия $l$.



П. п о в е р х в о с т и относительно точки $O-$ множество оснований перпевдикуляров, опущенных из точки $O$ на касательные плоскости поверхности. Уравнение П. поверхности $F(x, y, z)=0$ по отношению к началу координат:



\[

X=F_{x} \Phi, Y=F_{y} \Phi, Z=F_{z} \Phi,

\]



где



\[

\Phi=\frac{x F_{x}+y F_{y}+z F}{F_{x}^{2}+F_{y}^{2}+F_{z}^{2}}

\]



ПОДКАСАТЕЛЬНАЯ И ПОДНОРМАЛЬ - направленные отрезки $Q T$ и $Q N$, являющиеся проекциями на ось $O x$ отрезков касательной $M T$ и нормали $M N$ к век-рой кривой в ее точке $M$ (см. рис.).



Если кривая есть график функции $y=f(x)$, то значения величин П. и п. равны соответственно:



\[

Q T=-\frac{f(x)}{f^{\prime}(x)}, Q N=f(x) f^{\prime}(x),

\]



где $x-$ абсцисса точки $M$. Если кривая задана параметрически:



\[

x=\varphi(t), y=\psi(t),

\]



то тогда



\[

Q T=-\frac{\psi(t) \varphi^{\prime}(t)}{\psi^{\prime}(t)}, Q N=\frac{\psi(t) \psi^{\prime}(t)}{\psi^{\prime}(t)},

\]



где $t$ - значение параметра, определяющее точку $M$ кривой.



ПОДКАТЕГОРИЯ - частный случай понятия подструктуры математич. структуры. Категория $\mathfrak{R}$ наз. по дк а т е го р и ей категории $\Re, \quad$ если $\mathrm{Ob} \subseteq$ $\subseteq \mathrm{Ob} R$,



\[

H_{\mathfrak{R}}(A, B)=H_{\Re}(A, B) \cap \text { Mor } \mathfrak{R}

\]



для любых $A, B \in \mathrm{Ob} \mathfrak{\Omega}$ и произ ведение морфизмов из \& совпадает с их произведением в $\Re$. Для каждого





\end{document}